\chapter{Symétrie centrale et parallélogrammes}\label{ch:g7:central}

Tourner une figure d'un demi-tour autour d'un point~: c'est la
\emph{\omterm{def:g7:central:def}{symétrie centrale}}, la deuxième transformation du livre après
les réflexions du \cref{ch:g6:symmetry}. Sa figure vedette est le
\omterm{def:g7:central:parallelogram}{parallélogramme} --- le \omterm{def:g4:geometry:polygon}{quadrilatère} \omterm{def:g7:central:def}{symétrique} par rapport au point
de croisement de ses propres diagonales.

\section{Symétrie centrale}

\begin{definition}[Symétrique par rapport à un point]\label{def:g7:central:def}
Le \emph{symétrique}\index{symétrie centrale} d'un point $M$ par
rapport à un point $O$ est le point $M'$ tel que $O$ est le
\emph{milieu} de $[MM']$. La figure symétrique est ce que devient
l'original après un demi-tour ($180^\circ$) autour de $O$.
\end{definition}

\begin{omfigure}
\begin{tikzpicture}[scale=1.0]
  \fill (0,0) circle (1.6pt) node[below, font=\small] {$O$};
  \fill[omDef] (-2.3,1.0) circle (1.8pt) node[above, font=\small] {$M$};
  \fill[omThm] (2.3,-1.0) circle (1.8pt) node[below, font=\small] {$M'$};
  \draw[dashed] (-2.3,1.0) -- (2.3,-1.0);
  \draw (-1.2,0.45) -- (-1.1,0.6);
  \draw (1.1,-0.6) -- (1.2,-0.45);
  \fill[omDef] (-1.4,-0.8) circle (1.8pt) node[below, font=\small] {$N$};
  \fill[omThm] (1.4,0.8) circle (1.8pt) node[above, font=\small] {$N'$};
  \draw[dashed] (-1.4,-0.8) -- (1.4,0.8);
\end{tikzpicture}

{\small Symétrie par rapport à $O$~: chaque point et son image sont
alignés avec $O$, à distances égales de part et d'autre (traits de
coche).}
\end{omfigure}

\begin{method}[Construire le symétrique d'un point]\label{met:g7:central:construct}
Pour construire le \omterm{def:g7:central:def}{symétrique} $M'$ de $M$ par rapport à $O$~:
\begin{enumerate}
  \item tracer la droite $(MO)$ et la prolonger au-delà de $O$~;
  \item mesurer $MO$ (règle ou compas)~;
  \item placer $M'$ sur le prolongement avec $OM' = OM$.
\end{enumerate}
Sur papier quadrillé~: compter les pas horizontaux et verticaux de
$M$ à $O$, et répéter les \emph{mêmes} pas depuis $O$ pour atteindre
$M'$.
\end{method}

\begin{example}[Sur un quadrillage]\label{ex:g7:central:grid}
Si $M$ est à $3$ carreaux à gauche et $1$ carreau en haut de $O$, son
\omterm{def:g7:central:def}{symétrique} $M'$ est à $3$ carreaux à \emph{droite} et $1$ carreau en
\emph{bas} de $O$~: la \omterm{def:g7:central:def}{symétrie centrale} inverse les deux directions
à la fois (une réflexion n'en inverse qu'une).
\end{example}

\begin{proposition}[Ce qu'un demi-tour conserve]\label{prop:g7:central:preserved}
La \omterm{def:g7:central:def}{symétrie centrale} conserve les longueurs, les angles, les
\omterm{def:g6:measure:perimeter}{périmètres} et les \omterm{def:g6:measure:area}{aires}~; l'image d'une droite est une droite
\emph{\omterm{def:g6:lines:perp}{parallèle}}, l'image d'un \omterm{def:g6:lines:objects}{segment} est un \omterm{def:g6:lines:objects}{segment} \omterm{def:g6:lines:perp}{parallèle} de
même longueur, l'image d'un cercle est un cercle de même \omterm{def:g3:shapes:shapes}{rayon} (centré
au \omterm{def:g7:central:def}{symétrique} du centre).
\end{proposition}
\admitted

\begin{remark}
Comparer avec les réflexions (\cref{prop:g6:symmetry:preserved})~: les
deux conservent longueurs et angles~; mais la réflexion retourne les
figures (un «~b~» devient un «~d~»), tandis que le demi-tour les
garde lisibles --- à l'envers (un «~b~» devient un «~q~»). Et à la
\omterm{ex:g1:subtraction:difference}{différence} d'une réflexion, un demi-tour envoie toute droite sur une
droite qui lui est \omterm{def:g6:lines:perp}{parallèle}.
\end{remark}

\begin{definition}[Centre de symétrie]\label{def:g7:central:center}
Un point $O$ est un \emph{centre de
symétrie}\index{centre de symétrie} d'une figure lorsque le demi-tour
autour de $O$ envoie la figure exactement sur elle-même. Par exemple,
les lettres S, N, Z ont un centre de symétrie~; un cercle en a un (son
centre)~; un triangle n'en a jamais.
\end{definition}

\section{Parallélogrammes}

\begin{definition}[Parallélogramme]\label{def:g7:central:parallelogram}
Un \emph{parallélogramme}\index{parallélogramme} est un \omterm{def:g4:geometry:polygon}{quadrilatère}
dont les diagonales se croisent en leur milieu commun. Ce point de
croisement est alors un \omterm{def:g7:central:center}{centre de symétrie} de la figure~: chaque
\omterm{def:g5:solids:def}{sommet} est l'image par demi-tour du \omterm{def:g5:solids:def}{sommet} opposé.
\end{definition}

\begin{theorem}[Propriétés d'un parallélogramme]\label{thm:g7:central:properties}
Dans un \omterm{def:g7:central:parallelogram}{parallélogramme} $ABCD$ de centre $O$~:
\begin{enumerate}
  \item les côtés opposés sont \omterm{def:g6:lines:perp}{parallèles}~: $(AB) \parallel (CD)$ et
        $(BC) \parallel (AD)$~;
  \item les côtés opposés ont la même longueur~: $AB = CD$ et
        $BC = AD$~;
  \item les angles opposés sont égaux~: $\widehat A = \widehat C$ et
        $\widehat B = \widehat D$.
\end{enumerate}
\end{theorem}

\begin{proof}
Le demi-tour autour de $O$ envoie $A$ sur $C$ et $B$ sur $D$
(définition~: $O$ est le milieu des deux diagonales). Il envoie donc
le \omterm{def:g6:lines:objects}{segment} $[AB]$ sur le \omterm{def:g6:lines:objects}{segment} $[CD]$~; par
\cref{prop:g7:central:preserved} l'image est \omterm{def:g6:lines:perp}{parallèle} à $[AB]$ et de
même longueur~: points 1 et 2. Il envoie aussi l'angle en $A$ sur
l'angle en $C$, et les angles sont conservés~: point 3.
\end{proof}

\begin{omfigure}
\begin{tikzpicture}[scale=1.0]
  \coordinate (A) at (0,0);
  \coordinate (B) at (3.4,0);
  \coordinate (C) at (4.6,1.9);
  \coordinate (D) at (1.2,1.9);
  \draw[thick] (A) -- (B) -- (C) -- (D) -- cycle;
  \draw[omThm] (A) -- (C);
  \draw[omThm] (B) -- (D);
  \fill (2.3,0.95) circle (1.6pt) node[right=3pt, font=\small] {$O$};
  \node[below left, font=\small] at (A) {$A$};
  \node[below right, font=\small] at (B) {$B$};
  \node[above right, font=\small] at (C) {$C$};
  \node[above left, font=\small] at (D) {$D$};
  \draw (1.62,-0.09) -- (1.78,0.09);
  \draw (2.82,1.81) -- (2.98,1.99);
  \draw (0.53,1.02) -- (0.67,0.88);
  \draw (3.93,1.02) -- (4.07,0.88);
\end{tikzpicture}

{\small Un \omterm{def:g7:central:parallelogram}{parallélogramme} et son centre $O$~: les diagonales (rouges)
se coupent en leurs milieux, les côtés opposés sont \omterm{def:g6:lines:perp}{parallèles} et
égaux (traits de coche).}
\end{omfigure}

\begin{theorem}[Reconnaître un parallélogramme]\label{thm:g7:central:recognize}
Un \omterm{def:g4:geometry:polygon}{quadrilatère} (sans côtés croisés) est un \omterm{def:g7:central:parallelogram}{parallélogramme} dès que
\emph{l'une} des conditions suivantes est remplie~:
\begin{enumerate}
  \item ses diagonales ont le même milieu (la définition)~;
  \item ses côtés opposés sont \omterm{def:g6:lines:perp}{parallèles} deux à deux~;
  \item deux côtés opposés sont \omterm{def:g6:lines:perp}{parallèles} \emph{et} de même
        longueur.
\end{enumerate}
\end{theorem}
\admitted

\begin{example}\label{ex:g7:central:recognize}
$ABCD$ a $(AB) \parallel (CD)$ et $AB = CD = 5$ cm. Le critère 3
s'applique~: $ABCD$ est un \omterm{def:g7:central:parallelogram}{parallélogramme} --- donc sans autre mesure
on sait aussi que $AD = BC$, $(AD) \parallel (BC)$, et que ses
diagonales se coupent en leur milieu.
\end{example}

\begin{remark}[Parallélogrammes particuliers]
Le rectangle, le \omterm{def:g6:shapes:quadrilaterals}{losange} et le carré du \cref{ch:g6:shapes} sont
exactement les \omterm{def:g7:central:parallelogram}{parallélogrammes} avec une propriété
supplémentaire~: diagonales égales (rectangle), diagonales
\omterm{def:g6:lines:perp}{perpendiculaires} (\omterm{def:g6:shapes:quadrilaterals}{losange}), les deux (carré). L'argument d'angles
promis en \cref{ex:g6:shapes:recognize} fonctionne maintenant~: dans
un \omterm{def:g7:central:parallelogram}{parallélogramme}, les angles consécutifs sont supplémentaires
(\omterm{def:g7:triangles:pairs}{angles alternes-internes}/correspondants avec les \omterm{def:g6:lines:perp}{parallèles},
\cref{thm:g7:triangles:equalities}), donc un \omterm{def:g3:shapes:rightangle}{angle droit} force les
quatre.
\end{remark}

\section{Exercices}

\begin{exercise}[$\star$]\label{exo:g7:central:1}
Sur papier quadrillé, placer $O$, puis $M$ à quatre carreaux à droite
de $O$ et un carreau en haut. Construire le \omterm{def:g7:central:def}{symétrique} $M'$ de $M$
par rapport à $O$, et le \omterm{def:g7:central:def}{symétrique} de $M$ par rapport à la
\emph{droite verticale du quadrillage} passant par $O$. Comparer les
deux images.
\end{exercise}

\begin{exercise}[$\star$]\label{exo:g7:central:2}
Tracer un triangle $ABC$ et un point $O$ à l'extérieur. Construire
son image $A'B'C'$ par le demi-tour autour de $O$. Que peut-on dire
des longueurs $A'B'$ et $AB$~? Des droites $(A'B')$ et $(AB)$~?
\end{exercise}

\begin{exercise}[$\star$]\label{exo:g7:central:3}
Quels chiffres, écrits comme sur un afficheur de calculatrice ($0$ à
$9$), ont un \omterm{def:g7:central:center}{centre de symétrie}~? Quelles lettres majuscules parmi H,
A, N, S, T, Z, O~?
\end{exercise}

\begin{exercise}[$\star$]\label{exo:g7:central:4}
$KLMN$ est un \omterm{def:g7:central:parallelogram}{parallélogramme} de centre $O$ avec $KL = 7$ cm,
$LM = 4$ cm et $\widehat K = 115^\circ$. Donner, avec raisons~: $MN$,
$NK$, $\widehat M$, et le milieu de $[LN]$.
\end{exercise}

\begin{exercise}[$\star$]\label{exo:g7:central:5}
Construire un \omterm{def:g7:central:parallelogram}{parallélogramme} $ABCD$ à partir de ses diagonales~:
$AC = 8$ cm et $BD = 5$ cm, se croisant en leurs milieux avec un
angle de $60^\circ$ entre elles. (Tracer d'abord les diagonales~!)
\end{exercise}

\begin{exercise}[$\star$]\label{exo:g7:central:6}
Dans un \omterm{def:g7:central:parallelogram}{parallélogramme}, un angle mesure $115^\circ$. En utilisant la
remarque sur les angles consécutifs supplémentaires, donner les trois
autres angles.
\end{exercise}

\begin{exercise}[$\star\star$]\label{exo:g7:central:7}
Placer $A(1, 1)$, $B(4, 2)$ et le point $O(2.5, 2.5)$. Calculer les
coordonnées des images $A'$ et $B'$ par le demi-tour autour de $O$
(utiliser le comptage sur le quadrillage de
\cref{ex:g7:central:grid}).
\end{exercise}

\begin{exercise}[$\star\star$]\label{exo:g7:central:8}
$ABCD$ est un \omterm{def:g4:geometry:polygon}{quadrilatère} dans lequel $AB = CD$. Est-ce suffisant
pour en faire un \omterm{def:g7:central:parallelogram}{parallélogramme}~? Sinon, esquisser un
contre-exemple (un trapèze \omterm{def:g6:shapes:triangles}{isocèle} aide), et indiquer ce qu'il faut
ajouter à l'hypothèse.
\end{exercise}

\begin{exercise}[$\star\star$]\label{exo:g7:central:9}
Soit $ABC$ un triangle et $O$ le milieu de $[BC]$. Construire le
\omterm{def:g7:central:def}{symétrique} $A'$ de $A$ par rapport à $O$. Montrer que $ABA'C$ est un
\omterm{def:g7:central:parallelogram}{parallélogramme}. (Quel critère de \cref{thm:g7:central:recognize}
s'obtient gratuitement ici~?)
\end{exercise}

\begin{exercise}[$\star\star$]\label{exo:g7:central:10}
Vrai ou faux, avec une raison~: «~un \omterm{def:g4:geometry:polygon}{quadrilatère} dont les angles
opposés sont égaux deux à deux est un \omterm{def:g7:central:parallelogram}{parallélogramme}~»~; «~un
\omterm{def:g7:central:parallelogram}{parallélogramme} à diagonales \omterm{def:g6:lines:perp}{perpendiculaires} est un \omterm{def:g6:shapes:quadrilaterals}{losange}~»~; «~un
\omterm{def:g7:central:parallelogram}{parallélogramme} avec un \omterm{def:g3:shapes:rightangle}{angle droit} est un rectangle~».
\end{exercise}

\begin{exercise}[$\star\star\star$]\label{exo:g7:central:11}
Soit $ABCD$ un \omterm{def:g4:geometry:polygon}{quadrilatère} quelconque et $I$, $J$, $K$, $L$ les
milieux de ses côtés $[AB]$, $[BC]$, $[CD]$, $[DA]$. Tracer plusieurs
exemples. Que semble toujours être $IJKL$~? (La preuve, avec le
théorème des milieux, arrive au \cref{ch:g8:midpoints} --- et à
nouveau avec les vecteurs, dans un \omterm{def:g6:measure:volume}{volume} ultérieur.)
\end{exercise}

\section{Problème~: le gâteau, la médiane et l'impossible second
centre}

\begin{problem}[{Devoir maison --- le demi-tour à l'œuvre~: des
partages de gâteau équitables, un encadrement des médianes, et
pourquoi aucune figure ne peut avoir deux centres de symétrie}]
\label{pb:g7:central:1}
Le demi-tour a l'air de la plus simple des transformations --- et
pourtant il découpe les gâteaux en \omterm{def:g3:division:half}{moitiés} démontrablement
équitables, mesure les médianes d'un triangle, et cache une petite
perle de mathématiques pures~: une figure bornée peut avoir un \omterm{def:g7:central:center}{centre
de symétrie}, mais \emph{jamais deux}. La démonstration repose sur une
découverte que l'on fera dans la partie I~: deux demi-tours,
effectués l'un après l'autre, se ramènent à un glissement.

\medskip
\noindent\textbf{Partie I --- Gymnastique du demi-tour.}
\begin{enumerate}
  \item Sur du papier quadrillé, prendre l'origine $O(0, 0)$ pour
        centre. En comptant les carreaux comme dans
        \cref{ex:g7:central:grid}, donner les images de $M(3, 1)$,
        $N(-2, 4)$ et $P(0, -5)$ par le demi-tour de centre $O$.
        Quelle est la règle générale pour $(x, y)$~?
  \item Le centre est maintenant $C(2, 1)$. Calculer l'image de
        $M(5, 3)$, et vérifier la recette générale~: chaque
        coordonnée de l'image vaut \emph{le double de celle du
        centre, moins celle du point}.
  \item Les cartes à jouer sont dessinées pour se lire de la même
        façon à l'endroit et à l'envers --- une \omterm{def:g7:central:def}{symétrie centrale},
        de sorte qu'aucun joueur ne voie la carte «~de travers~».
        Expliquer pourquoi, sur une carte comptant un nombre
        \emph{\omterm{def:g3:numbers:evenodd}{impair}} de points, un point doit se trouver exactement
        au centre de la carte.
  \item Deux demi-tours d'affilée~: prendre les centres $O_1(2, 0)$
        et $O_2(5, 0)$. Envoyer le point $A(1, 1)$ par le demi-tour
        de centre $O_1$, puis envoyer l'image par le demi-tour de
        centre $O_2$. Comparer le départ et l'arrivée. Recommencer
        avec $B(0, 3)$. À quelle transformation unique et simple les
        deux demi-tours se sont-ils ramenés~?
  \item Comparer le glissement de la question 4 à la distance
        $O_1 O_2$, et énoncer la découverte. (Comparer avec les deux
        miroirs \omterm{def:g6:lines:perp}{parallèles} du~\cref{pb:g6:symmetry:1}~: même
        phénomène, nouveaux acteurs.)
\end{enumerate}

\medskip
\noindent\textbf{Partie II --- Gâteaux et médianes.}
\begin{enumerate}[resume]
  \item Soit $O$ le centre d'un \omterm{def:g7:central:parallelogram}{parallélogramme}. Expliquer pourquoi
        toute droite passant par $O$ rencontre le bord en deux
        points \omterm{def:g7:central:def}{symétriques} par rapport à $O$, et pourquoi cette
        droite découpe le \omterm{def:g7:central:parallelogram}{parallélogramme} en deux morceaux qui sont
        des copies exactes l'un de l'autre par demi-tour --- donc de
        même \omterm{def:g6:measure:area}{aire}.
  \item Le théorème du pâtissier~: un gâteau rectangulaire se
        partage parfaitement équitablement par \emph{n'importe
        quelle} coupe droite passant par son centre. Mieux~: un
        gâteau rectangulaire est percé d'un trou rectangulaire
        (n'importe quelle taille, n'importe quelle position, même
        incliné). Décrire l'unique coupe droite qui donne deux parts
        contenant autant de gâteau l'une que l'autre --- et la
        justifier par la question 6.
  \item Trois \omterm{def:g5:solids:def}{sommets} d'un \omterm{def:g7:central:parallelogram}{parallélogramme} $ABCD$ sont $A(0, 0)$,
        $B(6, 1)$ et $C(8, 5)$. Calculer les coordonnées de $D$ et
        celles du centre $O$.
  \item \Cref{exo:g7:central:9} a construit, à partir d'un triangle
        $ABC$ et du milieu $O$ de $[BC]$, le \omterm{def:g7:central:parallelogram}{parallélogramme} $ABA'C$
        où $A'$ est le \omterm{def:g7:central:def}{symétrique} de $A$ par rapport à $O$. S'en
        servir, avec l'inégalité triangulaire
        (\cref{thm:g7:triangles:inequality}) dans le triangle
        $ABA'$, pour démontrer l'\emph{encadrement de la médiane}~:
        la médiane $[AO]$ vérifie
        \[
        AO < \frac{AB + AC}{2} .
        \]
        (Remarquer que $AA' = 2\,AO$ et $BA' = AC$.)
  \item Dans un triangle où $AB = 5$ cm et $AC = 7$ cm, entre
        quelles deux valeurs la longueur de la médiane issue de $A$
        doit-elle se situer~? (Utiliser les deux sens de l'inégalité
        triangulaire dans $ABA'$.)
\end{enumerate}

\medskip
\noindent\textbf{Partie III --- Combien de centres une figure
peut-elle avoir~?}
\begin{enumerate}[resume]
  \item Pour chaque figure, dire si elle possède un \omterm{def:g7:central:center}{centre de
        symétrie}, et où~: un \omterm{def:g6:lines:objects}{segment}~; une droite entière~; un
        \omterm{def:g6:shapes:triangles}{triangle équilatéral}~; un carré~; un cercle~; les lettres S
        et Z.
  \item Démontrer la réponse pour le \omterm{def:g6:shapes:triangles}{triangle équilatéral}~: un
        demi-tour envoyant le triangle sur lui-même apparierait les
        trois \omterm{def:g5:solids:def}{sommets} --- mais trois est \omterm{def:g3:numbers:evenodd}{impair}, donc un \omterm{def:g5:solids:def}{sommet}
        serait envoyé sur lui-même. Qu'est-ce que cela obligerait ce
        \omterm{def:g5:solids:def}{sommet} à être, et pourquoi est-ce impossible~?
  \item Généraliser la question 12~: expliquer pourquoi
        \emph{aucun} \omterm{def:g4:geometry:polygon}{polygone} ayant un \omterm{def:g3:numbers:evenodd}{nombre impair} de \omterm{def:g5:solids:def}{sommets} n'a
        de \omterm{def:g7:central:center}{centre de symétrie}.
  \item La perle~: supposons qu'une figure possède \emph{deux}
        \omterm{def:g7:central:center}{centres de symétrie} distincts $O_1$ et $O_2$. Effectuer les
        deux demi-tours l'un après l'autre et utiliser la découverte
        de la question 5~: quel mouvement devrait envoyer la figure
        exactement sur elle-même~? Expliquer pourquoi une figure
        tenant sur une feuille de papier ne peut pas y survivre, et
        conclure.
  \item La conclusion de la question 14 laisse une échappatoire~:
        les figures \emph{non bornées}. Vérifier sur une frise
        infinie --- disons une bande sans fin d'empreintes de pas,
        gauche, droite, gauche, droite\,\dots\ --- qu'elle possède
        (au moins) deux \omterm{def:g7:central:center}{centres de symétrie} différents, et
        identifier la translation annoncée par la question 14.
\end{enumerate}
\end{problem}
